\documentclass[11pt]{article}
\usepackage[letterpaper,left=.9in,right=.9in,top=.8in,bottom=.85in]{geometry}
\usepackage{amsmath,amssymb,amsthm,xcolor,hyperref}
\hypersetup{colorlinks=true,linkcolor=blue!45!black,urlcolor=blue!45!black,citecolor=blue!45!black,
pdftitle={Irredundant covering systems with arbitrarily negative coprime-triple estimates},pdfauthor={GitHoobar}}
\newtheorem{theorem}{Theorem}
\title{\vspace{-1.5em}Irredundant covering systems with arbitrarily\\ negative coprime-triple estimates}
\author{GitHoobar}
\date{27 November 2026}
\begin{document}
\maketitle
\begin{abstract}
Majchrzak constructed full coverings for which the coprime-triple expression of
Harrington, Klein, Lowrance and Trifonov tends to minus infinity, and asked
what happens when every class of the full covering is essential. We construct
finite irredundant coverings of the integers by distinct moduli with the same
divergence. For every class we give a private integer, so removing any one
class destroys the covering.
\end{abstract}

For a finite set \(D\) of distinct integers greater than one, define
\[
 F(D)=\sum_{d\in D}\frac1d
 -\sum_{\substack{\{d,e\}\subseteq D\\(d,e)=1}}\frac1{de}
 +\sum_{\substack{\{d,e,f\}\subseteq D\\d,e,f\ {\rm pairwise\ coprime}}}\frac1{def}.
\]
The three sums use unordered sets of sizes one, two and three, respectively.
Harrington, Klein, Lowrance and Trifonov proved an upper bound on covered
density under a restriction on the prime factors and asked whether the
restriction could be removed \cite[Theorem 1.9 and Problem 4]{hklt}.
Majchrzak gave counterexamples that are full coverings, and separately
irredundant counterexamples whose union has density less than one.
He explicitly left open the case of irredundant full coverings
\cite[p.~1]{maj}. A private integer of a residue class belongs to that class
alone. A finite full covering is irredundant precisely when every class
has a private integer.

\begin{theorem}
For every real number \(R\) there is a finite irredundant covering of
\(\mathbb Z\) by one residue class for each \(d\in D\), with \(F(D)<R\).
\end{theorem}

\begin{proof}
We adapt the binary classes and the two coprime blocks of
\cite[p.~1]{maj}. We leave one residue modulo five open and attach a
different repair class to each hole of the odd block. This gives the new
classes private integers while preserving private integers for the old ones.
Choose any nonempty finite set \(Q\) of primes strictly greater than five.
Set \(N=3\prod_{q\in Q}q\). Begin with the classes
\[
 B_q:\quad x\equiv0\pmod{3q}\qquad(q\in Q).
\]
List the residues modulo \(N\) that they miss as \(u_1,\ldots,u_h\), with no
repetitions. They miss every integer coprime to 3, so \(h\ge 2N/3\ge4\).
Use the following additional classes:
\[
\begin{array}{lll}
\text{class} &\text{modulus}&\text{conditions on }x\\ \hline
P_i\ (1\le i\le h)&2^i&x\equiv 2^{i-1}-1\pmod {2^i}\\
M_j\ (1\le j\le4)&5\cdot2^j&x\equiv-1\pmod{2^j},\quad x\equiv j-1\pmod5\\
C_i\ (1\le i\le h)&5N\cdot2^i&x\equiv-1\pmod{2^i},\quad
                  x\equiv4\pmod5,\quad x\equiv u_i\pmod N .
\end{array}
\]
Each row specifies one class per listed modulus by the Chinese remainder
theorem. The even moduli have odd parts \(1,5,5N\), respectively, and the
powers of two distinguish them within each row. The odd moduli \(3q\) are
also pairwise distinct.

\emph{Coverage.} The classes \(P_i\) cover every residue other than
\(-1\bmod 2^h\): if \(2^h\nmid x+1\), take the first \(i\le h\) with
\(2^i\nmid x+1\). On the remaining residue all four binary conditions of
the \(M_j\) hold. These classes cover four residues modulo five.
The remaining points have \(x\equiv4\bmod5\). Their residue modulo \(N\)
is either covered by some \(B_q\) or equals \(u_i\), in which case \(C_i\)
covers them.

\emph{Irredundancy.} In the following table a triple gives the residues of a
private integer modulo \(2^h,5,N\), in that order. These three moduli are
pairwise coprime, so each triple is realized by an integer.
\[
\begin{array}{c|ccc}
\text{class}&x\bmod 2^h&x\bmod5&x\bmod N\\\hline
P_i&2^{i-1}-1&4&u_h\\
M_j&-1&j-1&u_h\\
B_q&-1&4&3q\\
C_i&-1&4&u_i
\end{array}
\]
In the first row only \(P_i\) has the indicated first failure of divisibility
of \(x+1\) by a power of two. Only \(C_h\) could match its \(N\)-coordinate;
that class requires \(x\equiv-1\bmod 2^h\). In the second row the residue
modulo five selects \(M_j\) alone among the mixed classes and excludes
every \(C_i\). Both rows use a hole of the \(B_q\).
In the third row \(3r\mid 3q\) for \(r\in Q\) holds only for \(r=q\);
as \(3q\) is covered, its residue is absent from the list of holes.
In the last row \(u_i\) selects exactly one repair class. This verifies
all four kinds of private integer.

\emph{Value of the expression.} Let \(A\) consist of all the pure classes
\(P_i\) and mixed classes \(M_j\), and let \(a,b,c\) be the sums of the reciprocal
moduli in \(A\), in the \(B_q\), and in the \(C_i\), respectively. The
coprime pairs are exactly the pairs from \(A\) and \(B\). Indeed, all
members of \(A\cup C\) are even, all members of \(B\cup C\) are divisible
by three, and every modulus in \(A\) has prime factors only in \(\{2,5\}\).
There are no pairwise-coprime triples. Consequently
\[
 F(D)=a+b+c-ab=1-(a-1)(b-1)+c,
\]
where
\[
 a=1-2^{-h}+\frac{3}{16}\ge\frac98,\qquad
 b=\frac13\sum_{q\in Q}\frac1q,\qquad
 c=\frac{1-2^{-h}}{5N}<\frac15 .
\]
By divergence of the sum of prime reciprocals, \(b\) can be made
arbitrarily large even after omitting the primes \(2,3,5\). For \(b\ge1\)
we therefore have \(F(D)<6/5-(b-1)/8\), which tends to minus infinity.
\end{proof}

\smallskip
\noindent\textit{Provenance and verification.} Developed with OpenAI Codex under
the GitHub handle GitHoobar. A
\href{https://github.com/GitHoobar/open-math-agents/blob/research/erdos-7/research/erdos-7/irredundant/IrredundantMain.lean}{Lean formalization}
verifies the theorem for every real threshold. The question about removing
redundant classes is due to Majchrzak. The coverings above contain even
moduli. The Erd\H{o}s--Selfridge odd covering problem remains open.

\begin{thebibliography}{9}
\bibitem{hklt} J.~Harrington, J.~Klein, J.~Lowrance and O.~Trifonov,
\emph{Covering systems where the prime divisors of all moduli are only 2, 3, or 5},
May 2026. \href{https://arxiv.org/abs/2605.18644}{arXiv:2605.18644}.
\bibitem{maj} F.~Majchrzak, \emph{Complete coverings for which a coprime-triple
density bound tends to minus infinity}, revision 2, November 2026.
\href{https://doi.org/10.5281/zenodo.24408596}{doi:10.5281/zenodo.24408596}.
\end{thebibliography}
\end{document}
